\chapter{Pairs and Baskets}\label{ch:s1:pairs-and-baskets}

Two synthetic stocks had moved together for a year, their normalised prices a standard deviation of 2.4 per cent apart: they were planted as close
substitutes, and the \omterm{def:s1:pairs-and-baskets:distance}{distance method} picked them. The pair opened when their spread reached two standard deviations, closed, and opened again. Then, on the 111th day of trading, one of the companies was taken over at a premium. The spread jumped to 35\%, the short leg stopped trading, and the pair
never closed: it lost 24.6\% on each dollar per leg. \omterm{def:s1:pairs-and-baskets:pairs}{Pairs trading} is the oldest story in \omterm{def:s1:anatomy-of-a-stat-arb-book:statarb}{statistical arbitrage} and the easiest to explain: find two stocks
that move together, and bet on their return to each other when they drift apart. This chapter is about how the two stocks are found, why most methods
find the wrong ones, and what happens when the story stops being true. The build is \texttt{firm.pairsel}.

\section{Selection: distance, cointegration, copulas}

\begin{definition}[Pairs trading]\label{def:s1:pairs-and-baskets:pairs}
\emph{Pairs trading}\index{pairs trading} is the strategy of selecting two securities whose prices have moved together, and taking a long position in the
one that has fallen behind and a short position in the other when their prices diverge, to be closed when they converge.
\end{definition}

\begin{definition}[Distance method, formation period, trading period]\label{def:s1:pairs-and-baskets:distance}
The \emph{distance method}\index{distance method} selects the pairs of stocks whose normalised price paths (cumulative total return indices starting at
one) have the smallest sum of squared differences over a \emph{formation period}\index{formation period}, and trades them over the following
\emph{trading period}\index{trading period}, opening a pair when its spread exceeds a multiple of its formation-period standard deviation and closing it
when the prices cross.
\end{definition}

Gatev, Goetzmann and Rouwenhorst gave the \omterm{def:s1:pairs-and-baskets:distance}{distance method} its standard form: twelve months of formation, six of trading, open at two historical
standard deviations, close at the next crossing or at the end of the six months; the top 20 pairs. On daily US data from 1962 to 2002 it earned annualised
excess returns of up to 11\% on self-financing portfolios of pairs, typically more than conservative cost estimates. The cointegration method (Book~4,
chapter~20) replaces the distance with an Engle--Granger test: regress one log price on the other and test the residual for a unit root. The copula method
models the pair's daily returns with a copula (Book~4, chapter~15) and trades when one return is improbably high or low given the other. Krauss's survey
counts over a hundred papers in five families, and Rad, Low and Faff compared the three main ones on US equities from 1962 to 2014: mean monthly excess
returns of 91, 85 and 43 basis points before costs for distance, cointegration and copula, and 38, 33 and 5 after.

The synthetic market so far has no pairs to find: its stocks share factors but no stock is a substitute for another, so every spread wanders off as the
difference of two random walks. \texttt{MarketConfig.twin\_share} (new in this chapter, off by default) pairs 10\% of the initial listings as close
substitutes, fifty pairs whose log price spread is a stationary AR(1) with a mean-reversion time of ten days and a standard deviation of 2\%. Everything
else about the twins is ordinary: they can be taken over, delisted or split like any stock. Because the pairs are known, the chapter can count what each
method finds.

\section{Trading rules and the formation window}

On the twin market, sixteen cycles of twelve months' formation and six months' trading, from year 3 to year 10, with the Gatev--Goetzmann--Rouwenhorst rule,
\$1 per leg and five basis points per unit traded:

\begin{center}
\small
\begin{tabular}{@{}lccccc@{}}
\toprule
selection & distance & copula & two-step & Engle--Granger & \omterm{def:s1:pairs-and-baskets:twin}{synthetic twin}\\
\midrule
planted pairs among the 320 selected & 311 & 311 & 113 & 40 & ---\\
Sharpe ratio before costs & 5.34 & 4.39 & 1.16 & $-0.09$ & $-0.45$\\
Sharpe ratio after costs & 5.05 & 3.89 & 1.03 & $-0.17$ & $-0.49$\\
net return a year (\$1 per leg) & 12.2\% & 11.6\% & 5.9\% & $-1.4\%$ & $-2.6\%$\\
openings per pair and cycle & 2.11 & 4.28 & 2.24 & 2.02 & 0.87\\
pairs still open at the end & 51\% & 89\% & 75\% & 81\% & 61\%\\
\midrule
Sharpe ratio after costs, no planted pairs & $-0.43$ & $-0.50$ & $-0.11$ & $-0.47$ & $-0.35$\\
\bottomrule
\end{tabular}
\end{center}

The \omterm{def:s1:pairs-and-baskets:distance}{distance method} finds the substitutes: 311 of its 320 pairs are planted twins, and it nets a Sharpe ratio of 5.05. On the default market, where there
are none, the same method loses: the pairs it picks are those whose paths happened to stay close over one year, typically two low-volatility stocks, and
84\% of them are still open when their six months end, because nothing pulls them back. Gatev, Goetzmann and Rouwenhorst noticed the same selection
effect: on average 82\% of the stocks in their top twenty pairs were utilities. The copula index on the same pairs trades twice as often and earns a little
less (its mispricing index drifts, so pairs stay open). The table's worst column is the one that sounds most rigorous.

\begin{omfigure}
\begin{tikzpicture}
\begin{axis}[width=11cm,height=5.0cm,xlabel={\small day of the trading period},ylabel={\small spread; P\&L per \$1 leg},tick label style={font=\footnotesize},
  xmin=0,xmax=126,ymin=-0.3,ymax=0.4,grid=major,grid style={gray!25},table/col sep=comma,legend pos=north west,legend style={font=\scriptsize}]
\addplot[omDef,very thick] table[x=day,y=spread]{figdata/strategies-1/04-pairs-and-baskets/takeover.csv};
\addplot[omThm,thick] table[x=day,y=pnl]{figdata/strategies-1/04-pairs-and-baskets/takeover.csv};
\addplot[black!50,dashed] coordinates {(0,0.048) (126,0.048)};
\addplot[black!50,dashed] coordinates {(0,-0.048) (126,-0.048)};
\addplot[black!60,dotted,thick] coordinates {(111,-0.3) (111,0.4)};
\legend{spread,cumulative P\&L,$\pm 2$ formation s.d.}
\end{axis}
\end{tikzpicture}

\omcaption{The chapter's opening example: two planted substitutes on the synthetic market over their six-month \omterm{def:s1:pairs-and-baskets:distance}{trading period}. The spread (difference of
normalised prices) opens the pair twice; on day 111 (dotted) the stock the pair is short is taken over at a premium, and the pair never closes. Data:
\texttt{s1\_pairs.takeover}.}\label{fig:s1:pairs-and-baskets:takeover}
\end{omfigure}

\section{Why cointegration tests fail at scale}

Ranking pairs by Engle--Granger p-values should find the planted pairs, since they are the only cointegrated ones. It finds 40 of 320. The reason is
multiple testing (Book~4, chapter~12) with a test whose p-values are wrong. Each cycle tests every pair of stocks in the same industry, about 46\,600 pairs;
over sixteen cycles, 745\,465 tests. The p-values come from the Engle--Granger statistic's distribution for two independent Gaussian random walks, simulated
once (\texttt{eg\_null}). The synthetic market's returns have fat tails and volatility that clusters, like real ones, and under those conditions the test
rejects too often: on the default market, where no pair is cointegrated, 7.4\% of pairs have p-values below 5\%, and 0.32\% below 0.1\%, three times the
nominal rate (\cref{fig:s1:pairs-and-baskets:calib}). The smallest p-values are dominated by these false rejections, not by the planted pairs.

\begin{omfigure}
\begin{tikzpicture}
\begin{axis}[width=7.2cm,height=5.6cm,xmode=log,ymode=log,xlabel={\small nominal level},ylabel={\small share of pairs rejected},
  tick label style={font=\footnotesize},xmin=0.0007,xmax=0.15,ymin=0.0007,ymax=0.2,grid=major,grid style={gray!25},legend pos=north west,
  legend style={font=\scriptsize},table/col sep=comma]
\addplot[omThm,very thick,mark=square*] table[x=nominal,y=actual]{figdata/strategies-1/04-pairs-and-baskets/calibration.csv};
\addplot[black!50,dashed] coordinates {(0.0007,0.0007) (0.15,0.15)};
\legend{Engle--Granger p-values,a valid test}
\end{axis}
\end{tikzpicture}

\omcaption{Calibration of the Engle--Granger test with a Gaussian random-walk null on the synthetic market without planted pairs (742\,028 same-industry pair
tests): the share of pairs whose p-value falls below each level, against the level. Data: \texttt{s1\_pairs.calibration}.}\label{fig:s1:pairs-and-baskets:calib}
\end{omfigure}

On the twin market, the counts are sharper still. Of 58\,024 pairs rejected at 5\% (7.8\% of those tested), 530 are planted, 84.7\% of the 626 planted pairs
the cycles could test: the test has power. Benjamini and Hochberg's procedure, set to a false discovery rate of 5\%, keeps 1\,139 pairs, of which 32 are
planted: 97.2\% of its discoveries are false. The procedure is correct; its input is not, and a false-discovery control built on miscalibrated p-values
controls nothing. Persistence tells the same story from the other side: of the pairs rejected at 5\% in a formation year and still listed a year later,
7.9\% are rejected again over that next year, which is barely above the test's false-rejection rate. The planted pairs are rejected again 84.6\% of the time,
the others 7.2\%.

The two-step method, a common repair, preselects the 200 pairs of smallest distance and then ranks them by Engle--Granger p-value. It keeps 113 planted
pairs and nets 1.03: better than the test alone, worse than the distance alone, because within the shortlist the test still prefers its false rejections.

\section{Why pairs died, and what replaced them}

The public record says profits declined. Do and Faff confirmed a continuing downward trend in pairs-trading profitability, with the strategy still
performing strongly in prolonged turbulence, including the 2008 crisis; Rad, Low and Faff found that from 2009 the distance and cointegration methods found
far fewer trading opportunities, while the copula method's stayed stable. The synthetic market gives the mechanism, if not the history: the \omterm{def:s1:pairs-and-baskets:distance}{distance method}
earns exactly as much as the market contains close substitutes, and nothing else. A method that selects on the past co-movement of prices finds
substitutes where they exist and low-volatility coincidences where they do not.

What replaced the classic pair is the rest of this part of the book: residuals against many factors (chapter~3), where the ``partner'' is a factor model;
structural relative value (chapter~12), where the substitutes are two share classes, a holding company and its holdings, or a fund and its assets, and the
link is contractual; and books of thousands of small bets (chapter~1), where no single pair matters.

\section{Baskets: one stock against its synthetic twin}

\begin{definition}[Synthetic twin]\label{def:s1:pairs-and-baskets:twin}
A stock's \emph{synthetic twin}\index{synthetic twin} is a portfolio of other stocks, usually from its industry, whose returns track the stock's over a
\omterm{def:s1:pairs-and-baskets:distance}{formation period}, estimated by a regression of the stock's returns on theirs; the spread between the stock and its twin is traded like a pair's.
\end{definition}

The basket generalises the pair: instead of one partner, thirty peers, weighted by a ridge regression on the formation year's returns. On the synthetic
market it loses in both versions ($-0.49$ and $-0.35$). A regression on thirty peers fits the stock's past factor moves well, but its specific moves are its
own: they are independent of every peer's, so the spread with any basket is the stock's specific random walk, and a random walk does not revert. The basket
works where the stock and its peers share something beyond the factors (a common customer, a regulated price), which is the question to ask of any
\omterm{def:s1:pairs-and-baskets:twin}{synthetic twin} before trading it.

\begin{omfigure}
\begin{tikzpicture}
\begin{axis}[width=9.5cm,height=5.4cm,xlabel={\small year},ylabel={\small cumulative net return (\%, \$1 per leg)},tick label style={font=\footnotesize},
  xmin=2,xmax=10,grid=major,grid style={gray!25},legend pos=outer north east,legend style={font=\scriptsize},table/col sep=comma]
\addplot[omDef,very thick] table[x=year,y=distance]{figdata/strategies-1/04-pairs-and-baskets/books.csv};
\addplot[omDef!60,thick] table[x=year,y=copula]{figdata/strategies-1/04-pairs-and-baskets/books.csv};
\addplot[omThm!60,thick] table[x=year,y=twostep]{figdata/strategies-1/04-pairs-and-baskets/books.csv};
\addplot[omThm,thick] table[x=year,y=eg]{figdata/strategies-1/04-pairs-and-baskets/books.csv};
\addplot[black!50,thick,dashed] table[x=year,y=twin]{figdata/strategies-1/04-pairs-and-baskets/books.csv};
\legend{distance,copula,two-step,Engle--Granger,synthetic twin}
\end{axis}
\end{tikzpicture}

\omcaption{Cumulative net returns of the five selection methods on the synthetic market with fifty planted pairs of substitutes. Data:
\texttt{s1\_pairs.run}.}\label{fig:s1:pairs-and-baskets:books}
\end{omfigure}

\section{Strategy files}

\begin{strategyfile}{Distance-method pairs}\label{strat:s1:pairs-and-baskets:distance}
\sfield{Who pays you, and why} Flows that push one of two close substitutes away from the other: the book lends liquidity across the pair.
\sfield{Instruments and venues} Liquid stocks, long one and short the other; the same exchange hours.
\sfield{Signal} The spread of normalised prices against its formation standard deviation; open at two, close at the crossing.
\sfield{Sizing and execution} \$1 per leg per pair, twenty pairs; both legs traded together at the close.
\sfield{Costs} Four units traded per round trip; borrow on the short leg.
\sfield{How it dies} The pairs are coincidences, not substitutes; a takeover or other news breaks a real pair (the chapter's opening example); too many traders
in the same pairs.
\sfield{Horizon, capacity, infrastructure} Weeks to months; small capacity per pair; a daily price file and a formation job.
\sfield{Backtest honestly} The formation universe as it was; delisting and takeover returns in the legs; pairs held to the end of the \omterm{def:s1:pairs-and-baskets:distance}{trading period}.
\sfield{Sources} Gatev, Goetzmann and Rouwenhorst (2006): up to 11\% a year in excess returns, 1962--2002; Do and Faff (2010): a continuing decline; Rad,
Low and Faff (2016): 38 basis points a month after costs, 1962--2014.
\end{strategyfile}

\begin{strategyfile}{Cointegration pairs}\label{strat:s1:pairs-and-baskets:coint}
\sfield{Who pays you, and why} As for distance pairs, selected by a test of a stationary spread.
\sfield{Instruments and venues} Liquid stocks, usually in the same industry.
\sfield{Signal} The Engle--Granger residual of one log price on the other, in formation standard deviations.
\sfield{Sizing and execution} Legs in the ratio of the cointegrating coefficient; open at two, close at zero.
\sfield{Costs} As for distance pairs.
\sfield{How it dies} Miscalibrated p-values across tens of thousands of tests fill the book with false pairs; relationships that break.
\sfield{Horizon, capacity, infrastructure} Weeks to months; a test per candidate pair per formation.
\sfield{Backtest honestly} Count the tests; check the test's calibration on data with no cointegration; preselect by economics, not by p-value alone.
\sfield{Sources} Rad, Low and Faff (2016): 33 basis points a month after costs; fewer opportunities after 2009.
\end{strategyfile}

\begin{strategyfile}{Copula pairs}\label{strat:s1:pairs-and-baskets:copula}
\sfield{Who pays you, and why} As for distance pairs; the signal measures how improbable today's joint move is.
\sfield{Instruments and venues} As for distance pairs.
\sfield{Signal} A mispricing index: the running sum of $P(U_a \le u_a \mid U_b = u_b) - \tfrac12$ under a copula fitted in the \omterm{def:s1:pairs-and-baskets:distance}{formation period}.
\sfield{Sizing and execution} Open at an index of $\pm 0.5$, close when it changes sign.
\sfield{Costs} Trades more often than the distance rule.
\sfield{How it dies} As above; the index drifts, and pairs stay open.
\sfield{Horizon, capacity, infrastructure} Weeks; marginal and copula fits per pair.
\sfield{Backtest honestly} Marginals and copula fitted on formation data only.
\sfield{Sources} Rad, Low and Faff (2016): 43 basis points a month before costs and 5 after; trading opportunities stable after 2009.
\end{strategyfile}

\begin{strategyfile}{Stock against its synthetic basket}\label{strat:s1:pairs-and-baskets:basket}
\sfield{Who pays you, and why} As for pairs, with a basket of peers as the partner.
\sfield{Instruments and venues} A stock and a basket of its industry peers.
\sfield{Signal} The cumulative residual of the stock against its ridge-regression twin, in formation standard deviations.
\sfield{Sizing and execution} \$1 in the stock against \$1 in the basket.
\sfield{Costs} The basket's trades as well as the stock's.
\sfield{How it dies} The stock's specific moves have no counterpart in any basket; the spread is then a random walk.
\sfield{Horizon, capacity, infrastructure} Weeks; a regression per stock per formation.
\sfield{Backtest honestly} Weights fitted on formation data only; compare with the stock's own \omterm{def:s1:short-term-reversal:residual}{residual reversal} (chapter~3).
\sfield{Sources} This chapter's simulation only; no public performance record verified.
\end{strategyfile}

\section{Tutorial: the pair that never closed}\label{tut:s1:pairs-and-baskets:tut}

\begin{tutorial}
\textbf{Goal.} Select pairs by distance, by Engle--Granger tests and by both; trade them; count the false discoveries. \textbf{End state:} the table and
\cref{fig:s1:pairs-and-baskets:calib}.
\begin{tutsteps}
\item \textbf{The tests}: Engle--Granger vectorised over many pairs, p-values from a simulated null, Benjamini--Hochberg.
\omcode{code/firm/pairsel/firm_pairsel.py}{55}{86}{The Engle--Granger statistic over many pairs, its null, and false-discovery control.}\label{lst:s1:pairs-and-baskets:eg}
\item \textbf{The rule}: open at two formation standard deviations, close at the crossing, positions drifting with returns.
\omcode{code/firm/pairsel/firm_pairsel.py}{89}{116}{The distance rule on one pair.}\label{lst:s1:pairs-and-baskets:rule}
\item \textbf{Run} \texttt{run(share, method)} for the five methods on both markets, \texttt{discoveries}, \texttt{calibration}, \texttt{takeover} and
\texttt{fig\_pairs.py}.
\end{tutsteps}
\textbf{What to change next.} Replace the Gaussian null by a bootstrap of each pair's own residuals and recount; widen the planted spread to 5\% and watch the
\omterm{def:s1:pairs-and-baskets:distance}{distance method} lose its twins to low-volatility coincidences; close pairs whose leg is under a takeover offer.
\end{tutorial}

\section{Build: pair selection}\label{bld:s1:pairs-and-baskets:pairsel}

\begin{build}
\bfield{Purpose} Selection of pairs by distance, cointegration and copula, and their trading rule, with the multiple-testing accounting that selection needs.
\bfield{Interface} \texttt{normalise(P)}, \texttt{top\_pairs(N, k)}, \texttt{eg\_tau(y, x)}, \texttt{eg\_null(T, reps, seed)}, \texttt{pvalue(tau, null)},
\texttt{bh(p, q)}, \texttt{trade\_pair(ra, rb, spread, sd, k, cost)}, \texttt{copula\_fit(ra, rb)}, \texttt{copula\_h(fit, ra, rb)}, \texttt{twin(y, X, lam)}.
\bfield{Rules} Formation data only in selection and scaling; every test counted; takeover and delisting returns in the legs.
\bfield{Acceptance tests} \texttt{code/firm/pairsel/tests/}: distance finds planted near-copies; the null's 5\% quantile near MacKinnon's value, the test's power
on a stationary spread and its size on random walks; Benjamini--Hochberg on a textbook example; the rule's P\&L by hand; the copula's correlation and
conditional probabilities; ridge weights; the inverse normal.
\bfield{Stretch} A bootstrap null per pair; Johansen tests on baskets (Book~4's \texttt{firm.coint}); stop-losses and takeover filters.
\end{build}

\begin{omsources}
\item E.~Gatev, W.~N.~Goetzmann and K.~G.~Rouwenhorst, ``Pairs trading: performance of a relative-value arbitrage rule'', \emph{Review of Financial Studies}
19(3), 2006 (NBER working paper 7032, 1999).
\item B.~Do and R.~Faff, ``Does simple pairs trading still work?'', \emph{Financial Analysts Journal} 66(4), 2010.
\item C.~Krauss, ``Statistical arbitrage pairs trading strategies: review and outlook'', \emph{Journal of Economic Surveys} 31(2), 2017.
\item H.~Rad, R.~K.~Y.~Low and R.~Faff, ``The profitability of pairs trading strategies: distance, cointegration and copula methods'', \emph{Quantitative
Finance} 16(10), 2016.
\end{omsources}

\section{Exercises}

\begin{exercise}[$\star$]\label{exo:s1:pairs-and-baskets:1}
Ten industries of 45 stocks each: how many same-industry pairs are there, and how many would a valid 5\% test reject if none is cointegrated?
\end{exercise}

\begin{exercise}[$\star$]\label{exo:s1:pairs-and-baskets:2}
Benjamini--Hochberg at $q = 0.05$ with 1\,000 tests: what p-value must the 20th smallest have for twenty rejections?
\end{exercise}

\begin{exercise}[$\star$]\label{exo:s1:pairs-and-baskets:3}
A pair is opened and closed once at five basis points per unit traded, \$1 per leg. What does the round trip cost, in basis points of one leg?
\end{exercise}

\begin{exercise}[$\star\star$]\label{exo:s1:pairs-and-baskets:4}
The opening example's formation standard deviation is 2.4\%. At what spread did the pair open, and why did the takeover make it lose far more than that?
\end{exercise}

\begin{exercise}[$\star\star$]\label{exo:s1:pairs-and-baskets:5}
Why does the \omterm{def:s1:pairs-and-baskets:distance}{distance method} favour pairs of low-volatility stocks, and why does that matter when the market has no substitutes?
\end{exercise}

\begin{exercise}[$\star\star$]\label{exo:s1:pairs-and-baskets:6}
Explain how 97\% of Benjamini--Hochberg's discoveries can be false when the procedure is set to 5\%.
\end{exercise}

\begin{exercise}[$\star\star\star$]\label{exo:s1:pairs-and-baskets:7}
\emph{Coding.} Using \texttt{discoveries(0.1)}, compute the share of planted pairs among the pairs rejected again in the following year, and compare it with
their share among the formation-year rejections.
\end{exercise}

\begin{exercise}[$\star\star\star$]\label{exo:s1:pairs-and-baskets:8}
\emph{Find the flaw.} ``We tested all 50\,000 same-industry pairs, kept those significant at 1\% after a false-discovery correction, and backtested the
survivors: Sharpe ratio 3.''
\end{exercise}

\section{Problem: The Pair That Never Closed}

\begin{problem}[{Weekend problem --- selection is the strategy}]\label{pb:s1:pairs-and-baskets:1}
The chapter's five selection methods on \texttt{firm.synthmkt}, with and without planted substitutes.

\medskip\noindent\textbf{Part I --- The methods.}
\begin{enumerate}
\item Define \omterm{def:s1:pairs-and-baskets:pairs}{pairs trading} and the \omterm{def:s1:pairs-and-baskets:distance}{distance method}'s formation and \omterm{def:s1:pairs-and-baskets:distance}{trading periods}.
\item State Gatev, Goetzmann and Rouwenhorst's rule and their headline result.
\item How do the cointegration and copula methods select and signal?
\item What did Rad, Low and Faff find, before and after costs?
\end{enumerate}

\medskip\noindent\textbf{Part II --- The synthetic test.}
\begin{enumerate}[resume]
\item What does \texttt{twin\_share} plant, and why does the chapter need it?
\item How many planted pairs does each method select?
\item Give the net Sharpe ratios on both markets.
\item Describe the opening example and why the pair never closed.
\end{enumerate}

\medskip\noindent\textbf{Part III --- Testing at scale.}
\begin{enumerate}[resume]
\item How many tests does the chapter run, and how well calibrated are they?
\item What share of Benjamini--Hochberg's discoveries is false, and why?
\item What does the two-step method do, and why does it not beat distance?
\item Why does the \omterm{def:s1:pairs-and-baskets:twin}{synthetic twin} lose?
\end{enumerate}

\medskip\noindent\textbf{Part IV --- The verdict.}
\begin{enumerate}[resume]
\item State the \emph{named result}: the share of cointegrated pairs in the formation year that stay cointegrated in the trading year, and the pairs book's
Sharpe ratio after costs.
\item What does persistence reveal about true and false pairs?
\item What did Do and Faff find about the trend and about crises?
\item What replaced the classic pair?
\item How would you check a test's calibration before using it at scale?
\item What would you require of a pair before trading it?
\item How would you protect a pairs book against takeovers?
\item In one sentence: what does a pairs strategy need from its market?
\end{enumerate}
\end{problem}

\section{Interview questions}

\begin{interviewq}[$\star$ \iqroles{researcher, trader}]\label{iq:s1:pairs-and-baskets:1}
How would you choose pairs for a pairs-trading strategy?
\end{interviewq}

\begin{interviewq}[$\star\star$ \iqroles{researcher}]\label{iq:s1:pairs-and-baskets:2}
Correlation or cointegration for pairs: what is the difference, and which matters?
\end{interviewq}

\begin{interviewq}[$\star\star$ \iqroles{researcher}]\label{iq:s1:pairs-and-baskets:3}
You test 100\,000 pairs for cointegration. How do you control false discoveries, and what can still go wrong?
\end{interviewq}

\begin{interviewq}[$\star\star$ \iqroles{risk, trader}]\label{iq:s1:pairs-and-baskets:4}
One leg of an open pair receives a takeover bid. What do you do?
\end{interviewq}

\begin{interviewq}[$\star\star$ \iqroles{developer}]\label{iq:s1:pairs-and-baskets:5}
Compute the sum of squared distances between the normalised prices of all pairs of 3\,000 stocks efficiently.
\end{interviewq}

\begin{interviewq}[$\star\star\star$ \iqroles{researcher}]\label{iq:s1:pairs-and-baskets:6}
A spread follows an AR(1) with coefficient $\phi$ and stationary standard deviation $\sigma$. What is the expected time to revert halfway, and how does the
Dickey--Fuller test's power depend on $\sigma$?
\end{interviewq}
